\chapter{Nhất Quỷ Nhì Ma}
\chapterintro{Những trò nghịch có truyền thống}{Một lớp học tử tế hiếm khi thiếu vài nhân vật chuyên gây biến: chuyền giấy, kéo ghế, giả giọng lớp trưởng, hô báo động giả rồi chạy như chưa từng quen biết.}

\poem{Bài 36}{Bài 36 - Chuyền giấy bay như bướm}{Passing tiny notes is an old-school art form.}{Nhất quỷ nhì ma nhé\\Giấy chuyền bay nhẹ như mơ\\Đáng ra vào tay đứa bạn\\Lại rơi trúng bàn cô cơ}

\poem{Bài 37}{Bài 37 - Kéo ghế nghe một cái rầm}{One pulled chair can echo through an entire term.}{Nhất quỷ nhì ma nhé\\Kéo ghế đúng lúc bạn ngồi\\Một tiếng rầm vang cả lớp\\Tôi cười, rồi chạy mất thôi}

\poem{Bài 38}{Bài 38 - Vặn đồng hồ sớm năm phút}{Five stolen minutes can feel like genius until discovered.}{Nhất quỷ nhì ma nhé\\Có đứa chỉnh lệch đồng hồ\\Mong chuông reo sớm năm phút\\Ai ngờ cô biết từ lâu}

\poem{Bài 39}{Bài 39 - Giả giọng lớp trưởng gọi tên}{A fake monitor voice can scare half the class.}{Nhất quỷ nhì ma nhé\\Có đứa giả giọng lớp trưởng\\Cả đám im phăng phắc hết\\Xong mới biết bị lừa thương}

\poem{Bài 40}{Bài 40 - Báo động giả làm cả lớp chạy}{Chaos travels faster than homework.}{Nhất quỷ nhì ma nhé\\Hô cô đến cả lớp dòm\\Mười cái đầu quay một lượt\\Bài ai nấy cất cái rụp}
